\section{Graphs and graph neural networks}
\label{sec:bg:graphs}
\label{sec:bg:graphs:concepts}

The transactions of the market, read as a whole, form a graph between wallets and markets,
and one of the families of the thesis reads it with message passing.  This section fixes
the vocabulary of graphs, gives the general form of a graph neural network and of the two
operators used here, attention over the neighbours and a gated update, and describes how
the same machinery is applied to a graph whose edges arrive one at a time.

\subsection{Graphs, bipartite graphs, dynamic graphs}

A \emph{graph} $G = (V, E)$ is a set of nodes $V$ and a set of edges $E \subseteq V \times
V$; its \emph{adjacency matrix} $\mathbf{A} \in \{0, 1\}^{|V| \times |V|}$ has $A_{uv} = 1$
when $(u, v) \in E$, and the \emph{neighbourhood} $N(v) = \{u : (u, v) \in E\}$ of a node
is the set of nodes joined to it.  The \emph{degree} of a node is the size of its
neighbourhood; a \emph{connected component} is a maximal set of nodes between which a
path exists; the \emph{$k$-core} is the largest sub-graph in which every node has at least
$k$ neighbours inside it, which isolates the dense centre of a graph with a long tail of
sparse nodes.  The transactions of a prediction market define a \emph{bipartite} graph:
$V = \mathcal{W} \cup \mathcal{M}$, wallets on one side and markets on the other, with an
edge between a wallet and a market for every trade of the first on the second and no edge
joining two wallets or two markets; its adjacency is a rectangular \emph{biadjacency}
matrix $\mathbf{B} \in \mathbb{R}^{|\mathcal{W}| \times |\mathcal{M}|}$, whose entries can
carry the number of fills between the pair.  A graph whose edges arrive one at a time,
each with a timestamp, is a \emph{dynamic} graph \cite{kazemi2020survey,gravina2024dynamic};
it can be read as a sequence of \emph{snapshots}, the graph of each hour, say, or as a
\emph{stream of events}, and the second reading is the one that matches the block-by-block
arrival of the fills.

\subsection{Message passing}
\label{sec:bg:gnn}

A \emph{graph neural network} computes a vector representation $\mathbf{h}_v$ of every
node from its own features and from the representations of its neighbours
\cite{scarselli2009gnn,micheli2009nn4g,kipf2017gcn,bacciu2020gentle}.  The general form,
the \emph{message-passing neural network} \cite{gilmer2017mpnn}, is one round of three
operations (\cref{fig:bg:mp}):
\begin{equation}
  \mathbf{msg}_{u \to v} = M(\mathbf{h}_v, \mathbf{h}_u, \mathbf{a}_{uv}), \qquad
  \mathbf{s}_v = \operatorname{AGG}_{u \in N(v)} \mathbf{msg}_{u \to v}, \qquad
  \mathbf{h}'_v = U(\mathbf{h}_v, \mathbf{s}_v),
  \label{eq:mpnn}
\end{equation}
a \emph{message} computed for every edge from the two ends and the edge's attributes
$\mathbf{a}_{uv}$ by a learned function $M$ (an MLP), an \emph{aggregation} of the
messages a node receives by a function that does not depend on their order or their
number, a sum, a mean or a maximum, and an \emph{update} of the node's representation
from its old value and the aggregate by a learned function $U$.  The aggregation is what
makes the network a function of the graph rather than of a particular listing of it: the
nodes of a neighbourhood have no natural order, and a sum or a mean gives the same result
in any order (\emph{permutation invariance}).  Several rounds let information travel
across the graph, one hop per round; the parameters of $M$ and $U$ are shared by all nodes
and all rounds in the simplest designs, so that the network can be applied to a graph of
any size and to nodes it has never seen (\emph{inductive} learning), which is necessary
here, where thirteen hundred new wallets appear every hour.  Training is backpropagation
through the three operations: the gradient of the loss at a node flows back through its
update, through the aggregation to every message, and through every message to the
representations of its neighbours and to the edge attributes, so that a node's parameters
are shaped by the predictions made at its neighbours as well as at itself.

\begin{figure}[t]
\centering
\resizebox{\textwidth}{!}{\begin{tikzpicture}[>=Latex, x=1cm, y=1cm,
  wal/.style={circle, draw=vizInk, fill=vizA!25, minimum size=7mm, font=\scriptsize, inner sep=1pt},
  mkt/.style={rectangle, draw=vizInk, fill=vizB!25, minimum size=7mm, font=\scriptsize, inner sep=2pt, rounded corners=1pt},
  edge/.style={draw=vizInk, line width=0.7pt},
  box/.style={draw=vizInk, rounded corners=2pt, align=center, font=\footnotesize, inner sep=3pt, minimum height=9mm, text width=46mm, fill=white},
  op/.style={box, fill=vizB!12}, state/.style={box, fill=vizA!12},
  arr/.style={->, draw=vizInk, line width=0.8pt}, lab/.style={font=\scriptsize, text=vizInk, align=center}]
% the bipartite graph of one block
\node[wal] (w1) at (0,2.0)  {$w_1$};
\node[wal] (w2) at (0,0.8)  {$w_2$};
\node[wal] (w3) at (0,-0.4) {$w_3$};
\node[wal] (w4) at (0,-1.6) {$w_4$};
\node[mkt] (m1) at (3,1.4)  {$m_1$};
\node[mkt] (m2) at (3,-1.0) {$m_2$};
\draw[edge] (w1) -- node[lab, above, sloped] {$\mathbf{a}_{e_1}$} (m1);
\draw[edge] (w2) -- node[lab, below, sloped] {$\mathbf{a}_{e_2}$} (m1);
\draw[edge] (w3) -- (m1);
\draw[edge] (w3) -- (m2);
\draw[edge] (w4) -- (m2);
\node[lab, anchor=north] at (1.5,-2.3) {the fills of one block: wallets (circles),\\markets (squares), one edge per fill\\with its attributes $\mathbf{a}_e$};
% the update of one market, as a vertical chain
\node[state] (hm)  at (8.6,3.0)  {state of $m_1$ before the block, $\mathbf{h}_{m_1}$};
\node[op]    (msg) at (8.6,1.5)  {one message per incoming edge\\$\mathbf{msg}_e = M(\mathbf{h}_{w}, \mathbf{h}_{m_1}, \mathbf{a}_e)$ for $e \in N_t(m_1)$};
\node[op]    (agg) at (8.6,-0.2) {aggregation, invariant to the order of the edges\\$\mathbf{s} = \sum_{e \in N_t(m_1)} \alpha_e\, \mathbf{msg}_e$};
\node[op]    (upd) at (8.6,-1.9) {update (a GRU cell)\\$\mathbf{h}_{m_1} \leftarrow U(\mathbf{h}_{m_1}, \mathbf{s})$};
\node[state] (hm2) at (8.6,-3.4) {state of $m_1$ after the block, $\mathbf{h}'_{m_1}$};
\draw[arr] (hm) -- (msg); \draw[arr] (msg) -- (agg); \draw[arr] (agg) -- (upd); \draw[arr] (upd) -- (hm2);
\draw[arr, dashed] (m1.east) -- (msg.west |- m1.east);
\node[lab, anchor=west, text width=28mm] at (11.6,-1.9) {then the same for every wallet, from the markets it traded in the block};
\end{tikzpicture}
}
\caption{Message passing on a bipartite graph of one block.  For a market, every incoming
fill produces a message from the wallet's state, the market's state and the fill's
attributes; the messages are aggregated in an order-independent way and written into the
market's state by an update function; the same is then done for every wallet from the
markets it traded.}
\label{fig:bg:mp}
\end{figure}

\paragraph{Attention over the neighbours.}  A plain sum treats every neighbour alike.
The \emph{graph attention network} \cite{velickovic2018gat} weighs the messages by a
learned relevance: for every edge a score is computed from the two ends (and, here, the
edge), the scores of the edges incident to a node are normalised with a softmax, and the
aggregation is the weighted sum,
\begin{equation}
  \alpha_{uv} = \frac{\exp\bigl(\operatorname{LeakyReLU}(\mathbf{q}^{\top} [\mathbf{h}_u, \mathbf{h}_v, \mathbf{e}_{uv}])\bigr)}
  {\sum_{u' \in N(v)} \exp\bigl(\operatorname{LeakyReLU}(\mathbf{q}^{\top} [\mathbf{h}_{u'}, \mathbf{h}_v, \mathbf{e}_{u'v}])\bigr)},
  \qquad
  \mathbf{s}_v = \sum_{u \in N(v)} \alpha_{uv}\, \mathbf{msg}_{u \to v},
  \label{eq:gat}
\end{equation}
where $[\cdot]$ is concatenation, $\mathbf{q}$ a learned vector and $\mathbf{e}_{uv}$ an
encoding of the edge attributes.  This is the attention of \cref{eq:attention} restricted
to the neighbours of a node, with one head.  \emph{Gated graph networks} \cite{li2016ggnn}
first used a GRU as the update function $U$, so that a node's representation is a state
that each round of message passing updates by \cref{eq:gru}, with the aggregate as input;
this is the update used in the stateful model, where a market's state is updated by the
book (\cref{sec:models:ctdgnn}) and, in the graph variant, by the fills of the block
(\cref{sec:models:mp}).  The depth of message passing raises stability questions that
\textcite{gravina2023antisymmetric} and \textcite{ceni2025mpssm} address with
anti-symmetric and dissipative updates; one of these, the continuous-time anti-symmetric
update, is tested in \cref{sec:stateful:variants}.

\subsection{Dynamic graphs: memories and messages in time}
\label{sec:bg:dyn}

A transaction stream is a \emph{continuous-time dynamic graph}: a sequence of timed events,
each adding an edge between a wallet and a market.  The event-based models of this family
keep a \emph{memory} vector per node and update it at every event.  Temporal Graph
Networks \cite{rossi2020tgn} give the general recipe, of which the earlier event-based
models \cite{kumar2019jodie,trivedi2019dyrep,xu2020tgat} are special cases: when an edge
$(u, v, t, \mathbf{a})$ arrives, a message is computed for each end from the two memories,
the elapsed time since the node's last event and the attributes, $\mathbf{msg}_v =
M(\mathbf{s}_v, \mathbf{s}_u, \Delta t_v, \mathbf{a})$; the messages a node receives in
a batch of events are aggregated; the memory is updated by a recurrent cell, $\mathbf{s}_v
\leftarrow \operatorname{GRUCell}(\bar{\mathbf{msg}}_v, \mathbf{s}_v)$; and the prediction
for a node is read from its memory, possibly after one more round of message passing over
the recent neighbours.  Training follows the stream in time order with truncated
backpropagation, exactly as in \cref{sec:bg:tbptt}, the batch being the events of a block.
The model of \cref{ch:models} is in this family: market and wallet states are memories,
the order-book window is a message a market sends to itself, and the transactions are
messages between wallets and markets.

Three results shaped the design.  \textcite{gravina2024ctan} (CTAN) propose a
non-dissipative update for continuous-time dynamic graphs that propagates information over
long ranges.  \textcite{feldman2026navis} revisit node-affinity prediction on temporal
graphs and show that a strong heuristic, the recent interaction history, beats most
learned models, and that a model should contain the heuristic as a special case; this is
the origin of the direct path of \cref{sec:models:skip}.  And benchmarks of spatiotemporal
forecasting raise the same warning: per-node baselines are competitive and must be reported
\cite{martin2026audit}, and when the graph is not given it can be learned sparsely from the
series \cite{cini2023sparse}.
